\subsection{Spectral Theory of Operators}

% **Role:** infinite-dimensional linear algebra powered by analysis.
%
% Spectrum, resolvent, spectral radius, self-adjoint/normal operators, functional calculus preview.

For bounded operators, spectrum replaces the finite-dimensional eigenvalue list. The resolvent tracks invertibility, while the spectrum records every obstruction to it.

\begin{definition} \label{def:bounded-operator-spectrum-resolvent} \hlt{Spectrum and Resolvent}\\
Let $T\in\LL(\mathcal{X})$ be a bounded operator on a Banach space. The \hlt{resolvent set} is
\[
\rho(T)\defeq\{\lambda\in\C:\lambda I-T\text{ is bijective and }(\lambda I-T)^{-1}\in\LL(\mathcal{X})\}.
\]
The \hlt{spectrum} is $\sigma(T)\defeq\C\setminus\rho(T)$.
\end{definition}

\begin{definition} \label{def:spectral-radius-bounded-operator} \hlt{Spectral Radius}\\
The \hlt{spectral radius} of $T\in\LL(\mathcal{X})$ is
\[
r(T)\defeq\sup\{\abs{\lambda}:\lambda\in\sigma(T)\}.
\]
\end{definition}

\begin{theorem} \label{thm:spectrum-nonempty-compact-bounded-operator} \hlt{Spectrum of a Complex Bounded Operator}\\
If $\mathcal{X}$ is a nonzero complex Banach space and $T\in\LL(\mathcal{X})$, then $\sigma(T)$ is a nonempty compact subset of $\C$ and $\sigma(T)\subseteq\{\lambda:\abs{\lambda}\leq\normop{T}\}$.
\end{theorem}

\begin{theorem} \label{thm:spectral-radius-formula} \hlt{Spectral Radius Formula}\\
For $T\in\LL(\mathcal{X})$ on a complex Banach space,
\[
r(T)=\lim_{n\to\infty}\normop{T^n}^{1/n}.
\]
\end{theorem}

\begin{perspective} \label{per:hilbert-operator-spectral-theory} \hlt{Hilbert Operator Spectral Theory}\\
For bounded self-adjoint and normal operators on Hilbert spaces, spectral theory recovers a functional-calculus form of diagonalization. Compact self-adjoint operators have an orthonormal eigenbasis away from the kernel, giving the closest infinite-dimensional analogue of the finite spectral theorem.
\end{perspective}
